\documentclass[11pt]{article}
\usepackage[margin=.76in,top=.62in,bottom=.62in]{geometry}
\usepackage{amsmath,amssymb,amsthm,hyperref}
\renewcommand{\normalsize}{\fontsize{10}{12}\selectfont}
\DeclareMathSizes{10}{10}{8}{6}
\hypersetup{colorlinks=true,linkcolor=blue,urlcolor=blue,citecolor=blue}
\setlength{\emergencystretch}{2em}
\newtheorem{theorem}{Theorem}
\newtheorem{lemma}[theorem]{Lemma}
\newcommand{\R}{\mathbb R}
\newcommand{\Z}{\mathbb Z}
\begin{document}
\begin{center}
{\Large\bfseries A failure of complete monotonicity for a Mills-ratio quotient}\\[.5em]
Research proof and Lean formalization\\
October 26, 2026
\end{center}
\begin{abstract}
We refute Ferreira--Simon Conjecture 3 on quotients of derivatives of the
Gaussian Mills ratio. Already the quotient $-r^2/r'$ has a negative
fortieth derivative throughout some interval of positive inputs. The proof
reduces the derivative at zero to an integer polynomial in $\pi/2$ of
degree 21. An exact rational Bernstein certificate proves its negativity
using six-decimal bounds on $\pi$; continuity gives a regular positive-input
counterexample. The literal Gaussian tail, derivative recurrences, finite
polynomial certificate, and positive-interval conclusion are formalized in Lean.
\end{abstract}

\section{The conjecture and the counterexample}
The Gaussian Mills ratio is
\begin{equation}\label{mills}
 r(x)=e^{x^2/2}\int_x^\infty e^{-t^2/2}\,dt,\qquad x\in\R.
\end{equation}
A function $g$ is completely monotone on $(0,\infty)$ if it is smooth there
and $(-1)^k g^{(k)}(x)\ge0$ for every integer $k\ge0$ and every $x>0$.
Ferreira and Simon~\cite{FS}, Conjecture 3 in arXiv v2, p.~31, conjecture
that
\[
 x\longmapsto-\frac{(r^{(n)}(x))^2}{r^{(2n+1)}(x)}
\]
is completely monotone for every nonnegative integer $n$. We use $n=0$.
The differential equation $r'=xr-1$ gives
\begin{equation}\label{F}
 F(x)=-\frac{r(x)^2}{r'(x)}=\frac{r(x)^2}{1-xr(x)}.
\end{equation}

\begin{theorem}\label{counterexample}
There is a $\delta>0$ such that, for every $0<x<\delta$,
\[
 r'(x)<0\qquad\hbox{and}\qquad F^{(40)}(x)<0.
\]
Consequently $F$ is not completely monotone on $(0,\infty)$, and
Ferreira--Simon Conjecture 3 is false.
\end{theorem}
The denominator is nonzero at every input in this conclusion. The result
concerns the source's uncomposed quotient. Their Proposition 10 establishes
complete monotonicity after composition with $x\mapsto\sqrt{x}$.

\section{An exact derivative certificate}
The Gaussian integral and differentiation of~\eqref{mills} give
\begin{equation}\label{ode}
 r(0)^2=\pi/2,\qquad r'(x)=xr(x)-1.
\end{equation}
Write $c=r(0)$, $R_j=r^{(j)}(0)$, and $D_j=F^{(j)}(0)$. The first identity
in~\eqref{F} is regular near zero because $r'(0)=-1$. Differentiating
\eqref{ode} and $F-xrF=r^2$ gives the finite recurrences
\begin{align}
 R_0&=c,\quad R_1=-1,\quad R_{j+2}=(j+1)R_j,\qquad D_0=c^2,\label{Rrec}\\
 D_{n+1}&=\sum_{j=0}^{n+1}\binom{n+1}{j}R_jR_{n+1-j}
 +(n+1)\sum_{j=0}^{n}\binom nj R_jD_{n-j}.\label{Drec}
\end{align}
These are polynomial identities over $\Z[c]$. Iterating~\eqref{Drec} to
$n+1=40$ gives
\begin{equation}\label{P}
 D_{40}=P(c^2),\qquad P\in\Z[T],\qquad\deg P=21.
\end{equation}
For orientation, the leading coefficient is $40!$ and numerical evaluation gives
\[
 P(\pi/2)\approx-2.4628065653022778370\times10^{22}.
\]
We certify the sign by rational arithmetic rather than by this decimal.
Equations~\eqref{Rrec}--\eqref{Drec} specify every coefficient of $P$;
the accompanying data gives its complete ascending coefficient list and
the 41 intermediate derivative polynomials.

Set
\[
 a=\frac{392699}{250000}=1.570796,\qquad
 b=\frac{3141593}{2000000}=1.5707965,\qquad h=b-a.
\]
The bounds $3.141592<\pi<3.141593$ imply $a<\pi/2<b$. If
$P(T)=\sum_{j=0}^{21}A_jT^j$, define the rational numbers
\begin{equation}\label{Bernstein}
 V_i=h^i\sum_{j=i}^{21}\binom ji A_ja^{j-i},\qquad
 B_k=\sum_{i=0}^k V_i\frac{\binom ki}{\binom{21}i}
 \quad(0\le k\le21).
\end{equation}
The power-to-Bernstein identity is
\begin{equation}\label{identity}
 P(a+hu)=\sum_{k=0}^{21}B_k\binom{21}k u^k(1-u)^{21-k}.
\end{equation}
Exact evaluation of the finite rational expressions in~\eqref{Bernstein}
gives, for each of the 22 indices,
\begin{equation}\label{bounds}
 -2463\cdot10^{19}<B_k<-2460\cdot10^{19}<0.
\end{equation}
The retained certificate stores both $B_k$ and $C_k=B_k\binom{21}k$.
Lean checks the full identity~\eqref{identity} with the $C_k$ and checks
$C_k<0$ for every index. Since $\binom{21}k>0$, these sign checks also
give $B_k<0$. The stronger rounded bounds~\eqref{bounds} are independently
checked by exact rational arithmetic in the manuscript QA record.

Take $u=(\pi/2-a)/h\in(0,1)$. Every summand in~\eqref{identity} is
strictly negative, so $P(\pi/2)<0$. By~\eqref{ode} and~\eqref{P},
\[
 F^{(40)}(0)<0.
\]

To finish Theorem~\ref{counterexample}, the Gaussian tail is smooth on
$\R$, and the denominator of~\eqref{F} is 1 at zero. Thus $F^{(40)}$ is
continuous near zero. Its strict negative value persists on a
neighborhood, as does $r'(0)=-1<0$. Choose $\delta>0$ inside both
neighborhoods. For every $0<x<\delta$, the even-order complete-monotonicity
inequality fails at order 40. This proves the theorem on genuinely positive
inputs.

\section{Formalization, reproducibility, and scope}
The Lean certificate defines $r$ by the literal improper Gaussian integral
in~\eqref{mills}. It proves the finite-integral bridge, the initial value,
the differential equation, and smoothness. Finite Leibniz identities link
the 41 archived integer polynomials to the genuine iterated derivatives of
\eqref{F}; all table identities are proved. The Bernstein identity and
signs are then checked exactly. The endpoint
\begin{center}
\texttt{MillsCounterexample.source\_counterexample\_351688}
\end{center}
is a premise-free negation of the literal source conjecture. The stronger
interval endpoint also proves that $r'$ stays negative there. There is no
residual differential-equation, polynomial-identity, or numerical-bound
assumption in these endpoints.

The certificate has 47 declarations in seven modules, built against Lean
4.34.1 and mathlib \texttt{d13f23b7}. The replay procedure copies source
files afresh, checks pinned dependencies, treats warnings as errors,
inspects every declaration's type and axioms, and excludes proof holes,
extra axioms, and native decision procedures. The exact verification script
independently reconstructs the recurrences and Bernstein identity. A
separate Cauchy quadrature diagnostic of the complex error-function formula
agrees with the derivative value to relative error below $4\times10^{-83}$;
this numerical diagnostic is not used in the proof.

This is one refutation, already at $n=0$. We make no assertion about the
least derivative order of failure or the remaining individual values of
$n$. The source comparison uses the latest arXiv v2; the published full-text
access limit is recorded. The bounded prior-work investigation and
independent review remain pending for this draft, and no absolute priority
claim is made. Later complete-monotonicity results such as~\cite{Yang}
concern different weighted or composed ratios.

\begin{thebibliography}{9}
\bibitem{FS} R. A. C. Ferreira and T. Simon,
\emph{Convolution of beta prime distribution}, Trans. Amer. Math. Soc.
\textbf{376} (2023), 855--890.
\href{https://doi.org/10.1090/tran/8748}{doi:10.1090/tran/8748};
\href{https://arxiv.org/abs/2108.09244v2}{arXiv:2108.09244v2}.
\bibitem{Yang} Z.-H. Yang, \emph{Extension of a complete monotonicity theorem
with applications}, \href{https://arxiv.org/abs/2507.10954}{arXiv:2507.10954}
(2025).
\end{thebibliography}
\end{document}
